\section{The Mother of God and the Coming Year}

The Church teaches us to ask the Mother of God's intercession for the hour of death and to entrust ourselves to St.~Joseph, patron of a happy death (CCC 1014). These familiar devotions place the unknown hour within the communion of saints. We are asking for help at an event that will really come, for perseverance in grace when our own resources may be least available. The habitual petition becomes more deliberate as we allow its meaning to reach the life we are living.

Mary stood by the cross of her Son, where he gave her to the beloved disciple as his mother (Jn 19:25--27). She remained with him through his suffering. When we ask her help for our last hour, we seek her intercession with the Son whose death she witnessed. We ask to remain faithful to him, to receive his mercy, and to endure the trial with her help. Her presence at Calvary gives us reason to turn to her when the suffering of someone we love leaves us able only to remain beside him and pray.

We also ask St.~Joseph to help us persevere faithfully. The Church commends him as patron of a happy death; the Gospels recount his care for Jesus and Mary while leaving his own death unrecorded. We can turn to his intercession as we put our affairs in order and seek the grace to surrender them. The responsibilities we have received remain duties to fulfill, and we can ask his help in providing for those entrusted to us before our work passes to others.

In asking these saints' help, we seek the grace to repent, forgive, repair, and persevere. Their intercession assists the sacramental life and charity in which we are called to live. We can receive the prompting to confess a concealed fault, follow a priest's counsel, or undertake a resisted duty with gratitude for the help we have sought. A familiar devotion becomes more deliberate when we remember what we have asked for and respond to the opportunities for amendment that follow.

An annual return to the Four Last Things gives us time to examine the year that has passed and consider the year we may still receive. We can give thanks for the graces received in the sacraments, for protection and reconciliation, and for the people through whom we have been loved. We can acknowledge the sins that require confession and the duties still awaiting attention. Prayer can then address these needs specifically, asking for help with the work and relationships in which faithfulness will be required of us.

A quiet period for this reflection can begin with Christ at Lazarus's tomb, his prayer in Gethsemane, or his mercy toward the dying thief. Attending to what he does and says helps us bring our needs to him in faith. We can then examine the year. Has a recurring fault become easier to excuse? Has prayer been steadily crowded out? Have those who depend on us received patience and attention? Are we allowing an old injury to determine how we treat someone? We can consider a particular instance and ask what repentance or greater charity now requires.

The next step should be proportionate and definite. A confession can be arranged. A needed restitution can be begun. A conversation with a priest can address a matter conscience cannot resolve alone. One neglected duty can receive a place in the week. A resolution becomes more trustworthy when it takes bodily form in time, money, words, or service. The annual return then has a fruit that persists after the emotional force of the meditation has faded.

Practical affairs also belong within charity. The people who survive us may need to know where essential records are kept, what responsibilities remain, whom to contact, and what we desire for sacramental care and a Catholic funeral. Appropriate legal and medical arrangements require competent advice and attention to actual circumstances. Their spiritual purpose is plain: we are relieving others of avoidable confusion and making provision for those entrusted to us. Arranging these matters can be an act of humility, an acknowledgment that our work will pass into other hands.

The remembrance of our dead should have a place as well. Their names can be brought to Mass and prayer; an anniversary can renew intercession; the needs of someone they once cared for can be attended to. We may discover that grief has become closed around a single painful scene. The Church's prayer opens it again toward the whole person and toward the mercy of Christ. We can remember the gifts with gratitude, acknowledge what was difficult, and ask for the purification and fulfillment that only God can give.

The coming year remains uncertain. If it is given, there will be work to do and pleasures for which to give thanks. We can return to them with a clearer sense of the time available to love the people still with us. A visit we have postponed can be made; an apology can be offered; a kindness can be acknowledged. Some circumstances will remain unresolved despite our efforts. We can place them before God and continue doing the good within our power, grateful for the life in which he still permits us to serve.

